% BHU PYQ -- Manually transcribed & error-corrected
% Paper: BCH-216 : Basic Statistics
% Exam: B.Com. (Hons.) Semester III Examination, 2013-14

\documentclass[11pt,a4paper]{article}
\usepackage[a4paper,top=2.5cm,bottom=2.5cm,left=2.8cm,right=2.8cm,headheight=14pt]{geometry}
\usepackage{amsmath,amssymb}
\usepackage[T1]{fontenc}
\usepackage[utf8]{inputenc}
\usepackage{lmodern,microtype,fancyhdr,enumitem,hyperref,lastpage,parskip}

\hypersetup{
    colorlinks=true,
    urlcolor=black,
    linkcolor=black,
    pdftitle={BCH-216: Basic Statistics -- BHU B.Com. (Hons.) Sem III 2013-14}
}

\pagestyle{fancy}
\fancyhf{}
\renewcommand{\headrulewidth}{0.4pt}
\renewcommand{\footrulewidth}{0.4pt}
\fancyhead[L]{\small   \textit{Banaras Hindu University}}
\fancyhead[R]{\small   \textit{BCH-216: Basic Statistics}}
\fancyfoot[C]{\small Page  \thepage\ of \pageref{LastPage}}


\newlist{parts}{enumerate}{2}
\setlist[parts,1]{label=(\alph*),leftmargin=*,itemsep=5pt,topsep=3pt}
\setlist[parts,2]{label=(\roman*),leftmargin=*,itemsep=3pt,topsep=2pt}

\newcommand{\pts}[1]{\hfill   \textit{[#1]}}


\begin{document}


\begin{center}
  {\large\bfseries BANARAS HINDU UNIVERSITY}\\[2pt]
  {\normalsize B.Com. (Hons.) --- Semester III Examination, 2013-14}\\[6pt]
  \rule{0.8\linewidth}{0.5pt}\\[6pt]
  {\large\bfseries Paper No. BCH-216: Basic Statistics}\\[4pt]
  {\normalsize Time: 3 Hours\hfill Maximum Marks: 70}
\end{center}

\rule{\linewidth}{0.4pt}

\smallskip



\noindent   \textit{Instructions: Answer all questions. All questions carry equal marks (14 marks each).}

\bigskip




\noindent   \textbf{Question 1.} \\
What is secondary data? What methods are used to collect these? Explain their relative merits and demerits.\pts{14}

\centerline{   \textbf{OR}}

Frequency distribution of marks obtained by a class of students shows the following:\pts{14}

\begin{center}
\renewcommand{\arraystretch}{1.2}

\begin{tabular}{|l|c|c|c|c|c|c|}
\hline
   \textbf{Marks} & 0--30 & 30--40 & 40--50 & 50--60 & 60--70 & 70--100 \\ \hline
   \textbf{Frequency (f)} & 10 & 15 & 30 & 32 & 8 & 5 \\ \hline
\end{tabular}
\end{center}

\begin{parts}
  \item Find the median by drawing the ogive curve.
  \item Check up the value of median.
\end{parts}

\medskip\rule{\linewidth}{0.2pt}




\noindent   \textbf{Question 2.} \\
An aeroplane covered a distance of 800 kilometres with four different speeds of 100, 200, 300 and 400 kmph for the first, second, third and fourth quarter of a distance respectively. Find the average speed in kmph. Test the validity of your answer.\pts{14}

\centerline{   \textbf{OR}}

Find out mode and median wages of the following data:\pts{14}

\begin{center}
\renewcommand{\arraystretch}{1.2}

\begin{tabular}{|l|c|c|c|c|c|}
\hline
   \textbf{Wages (Rs.)} & 0--10 & 10--20 & 20--30 & 30--40 & 40--50 \\ \hline
   \textbf{No. of workers} & 5 & 7 & 15 & 25 & 8 \\ \hline
\end{tabular}
\end{center}

\medskip\rule{\linewidth}{0.2pt}




\noindent   \textbf{Question 3.} \\
Goals scored by teams A and B in a football season were as follows:\pts{14}

\begin{center}
\renewcommand{\arraystretch}{1.2}

\begin{tabular}{|c|c|c|}
\hline
   \textbf{No. of goals} & \multicolumn{2}{c|}{   \textbf{No. of matches}} \\ \cline{2-3}
&   \textbf{Team A} &   \textbf{Team B} \\ \hline
0 & 27 & 17 \\
1 & 9 & 9 \\
2 & 8 & 6 \\
3 & 5 & 5 \\
4 & 4 & 3 \\ \hline
\end{tabular}
\end{center}
By calculating the coefficient of variation in each, find which team may be considered more consistent.

\centerline{   \textbf{OR}}

If the values of mean, median and mode are Rs. 33.26, Rs. 33.5 and Rs. 34 respectively in the following distribution, calculate the missing frequencies and also Karl Pearson's coefficient of skewness:\pts{14}

\begin{center}
\renewcommand{\arraystretch}{1.2}

\begin{tabular}{|l|c|c|c|c|c|c|c|c|}
\hline
   \textbf{Wages (Rs.)} & 0--10 & 10--20 & 20--30 & 30--40 & 40--50 & 50--60 & 60--70 &   \textbf{Total} \\ \hline
   \textbf{Frequencies} & 8 & 32 & ? & ? & ? & 12 & 8 & 460 \\ \hline
\end{tabular}
\end{center}

\medskip\rule{\linewidth}{0.2pt}




\noindent   \textbf{Question 4.} \\
What is Spearman's rank correlation coefficient? Bring out its usefulness.\pts{14}

\centerline{   \textbf{OR}}

From the following data, find out the coefficient of correlation between age and illiteracy:\pts{14}

\begin{center}
\renewcommand{\arraystretch}{1.2}

\begin{tabular}{|l|c|c|}
\hline
   \textbf{Age group} &   \textbf{Total population} &   \textbf{Illiterate population} \\ \hline
10--20 & 120 & 100 \\
20--30 & 100 & 75 \\
30--40 & 80 & 60 \\
40--50 & 50 & 30 \\
50--60 & 25 & 20 \\
60--70 & 15 & 10 \\
70--80 & 5 & 5 \\ \hline
\end{tabular}
\end{center}

\medskip\rule{\linewidth}{0.2pt}




\noindent   \textbf{Question 5.} \\
Distinguish between correlation and regression. Also discuss the various methods of measuring regression.\pts{14}

\centerline{   \textbf{OR}}

From the following data, calculate:\pts{14}

\begin{parts}
  \item the two regression coefficients;
  \item the two regression equations;
  \item estimate the likely height of the son when father's height is 67.5 inches;
  \item estimate the likely height of father when son's height is 68.2 inches.
\end{parts}

\begin{center}
\renewcommand{\arraystretch}{1.2}

\begin{tabular}{|l|c|c|c|c|c|c|c|c|}
\hline
   \textbf{Height of father (inches)} & 65 & 66 & 67 & 67 & 68 & 69 & 71 & 73 \\ \hline
   \textbf{Height of son (inches)} & 67 & 68 & 64 & 68 & 72 & 70 & 69 & 70 \\ \hline
\end{tabular}
\end{center}

\vfill
\rule{\linewidth}{0.4pt}

\begin{center}\small
\href{https://bhu-exam.vercel.app/}{Click here to visit our website}\\[4pt]
\href{https://me-aryan.vercel.app/}{Click here to visit our contributor}\\[4pt]
\href{https://quashan-me.vercel.app/}{Click here to visit our contributor}
\end{center}

\end{document}
